
\begin{center}
  （与 Laurent Lazzarini 合写）
\end{center}

本附录说明关于 \BookRecordIndex{T0003}$J$-全纯曲线的两种相关现象。第一种涉及奇点附近局部结构的精细分析
（全文假设所论曲线没有边界）。这部分的主要结果断言：\BookRecordIndex{T0334}伪全纯曲线的奇点可由一个
$C^1$-微分同胚共轭到全纯曲线的奇点~\srcite{287}。因此，每条单纯
\BookRecordIndex{T0003}$J$-全纯曲线只有有限多个自交点，两条不同的 \BookRecordIndex{T0003}$J$-全纯曲线也只能在有限多个点相交。\sourcefootnote{原书导言此句未写\BookRecordIndex{T0192}紧性或像不共享开集等限定；下文定理 E.1.2 对\BookRecordIndex{T0119}非单射点有限性要求定义域闭，定理 E.1.5 对两曲线还要求不存在像相同的非空开片。此处保留导言表述，精确适用条件以下文定理为准。}
第二种现象为四维所特有：两条不同 \BookRecordIndex{T0003}$J$-全纯曲线的每个交点都对交数作正贡献，
每个奇点也都对自交数作正贡献。由此得到\BookRecordIndex{T0054}伴随不等式
\[
  2g-2\le C\cdot C-\langle c_1(TM),C\rangle,
\]
它对 $M^4$ 中每条单纯 \BookRecordIndex{T0003}$J$-全纯曲线 $C$ 成立，且等号成立当且仅当 $C$ 为嵌入曲线。
局部分析的另一推论是：凡可由 \BookRecordIndex{T0003}$J$-全纯曲线表示的同调类
$A\in H_2(M;\mathbb Z)$，也可由某个浸入的 $J'$-全纯曲线表示，其中 $J'$ 在
$C^1$ 意义下接近 $J$。这一结论在任意维度都成立。

这些结果都是已知的；本附录旨在给出完整的阐述。Gromov~\srcite{160} 已指出交点的
有限性与正性。McDuff~\srcite{256} 首次给出了证明（勘正见~\srcite{264}），
该文也是首先陈述并证明\BookRecordIndex{T0054}伴随不等式的文献。Micallef 与 White~\srcite{287} 从另一
途径研究了这些问题。他们处理任意\BookRecordIndex{T0175}极小曲面，再将所得结果用于 \BookRecordIndex{T0003}$J$-全纯曲线这一
特例，因此其工作更为一般。Sikorav~\srcite{379} 给出了这一系列思想的第三种阐述。
他利用 \BookRecordIndex{T0013}Carleman 相似原理，得到了把伪全纯奇点共轭到全纯奇点的微分同胚的尖锐
\BookRecordIndex{T0408}正则性结果：它可取为 $C^{1+\mu}$ 类，其中 $\mu<1$，但未必为 $C^2$ 类。
该文发表于下述证明接近完成之际。

本附录安排如下。首先陈述 \BookRecordIndex{T0036}Micallef--White 定理，以及关于交点有限性和正性的主要
结果。E.2 节说明这些结果如何从 \BookRecordIndex{T0036}Micallef--White 定理\BookRecordIndex{T0326}推出。后者的证明又建立在
\BookRecordIndex{T0024}Hartman--Wintner 定理之上，该定理描述\BookRecordIndex{T0194}近似调和函数的最低阶项。要在所需的一般性
下证明它，须用到关于二维\BookRecordIndex{T0088}殆复结构可积性的经典结果（\BookRecordIndex{T0041}Newlander--Nirenberg 定理）
及\BookRecordIndex{T0098}等温坐标的存在性。E.3 节确立这些结果后，E.4 节证明 \BookRecordIndex{T0024}Hartman--Wintner 定理；
这些内容都相当标准。本方法的新意在于用 \BookRecordIndex{T0024}Hartman--Wintner 定理理解两条相交的
\BookRecordIndex{T0003}$J$-全纯曲线之差 $u-v$ 的最低阶项。这一步占据 E.5、E.6 两节。最后的 E.7 节说明
Micallef--White 的论证：如何在局部将 \BookRecordIndex{T0003}$J$-全纯曲线共轭到复多项式。

\section{主要结果}\label{sec:E-1}

本节陈述本附录所证明的主要结果。全文中，$(M,J)$ 是维数为 $2n$ 的光滑殆复流形，
不要求紧。$M$ 中一条 \BookRecordIndex{T0003}$J$-全纯曲线是定义在带复结构的二维流形 $(\Sigma,j)$ 上的光滑映射
$u:\Sigma\to M$，不要求 $\Sigma$ 紧，并满足
\[
  du\circ j=J\circ du.
\]
若不存在两个不交的非空开集 $U_0,U_1\subset\Sigma$ 使
$u(U_0)=u(U_1)$，则称 $u$ 为\emph{单纯的}。若 $\Sigma$ 是闭的，这意味着 $u$
在 $\Sigma$ 的每个连通分支上的限制均非多重覆盖，且不同分支在 $u$ 下的像不同。
给定 \BookRecordIndex{T0003}$J$-全纯曲线 $u:\Sigma\to M$，若 $du(z)=0$，则称 $z\in\Sigma$ 为 $u$ 的
\emph{\BookRecordIndex{T0233}临界点}；若 $z$ 不是\BookRecordIndex{T0233}临界点且 $u^{-1}(u(z))=\{z\}$，则称它为 $u$ 的
\emph{\BookRecordIndex{T0094}单射点}。

回顾引理~\ref{lem:2-4-1}：定义在闭黎曼曲面 $\Sigma$ 上的每条
\BookRecordIndex{T0003}$J$-全纯曲线 $u:\Sigma\to M$ 都只有有限多个\BookRecordIndex{T0233}临界点，而且是有限对一的，
即对每个 $x\in M$，集合 $u^{-1}(x)$ 有限。此外，命题~\ref{prop:2-5-1} 断言，
单纯 \BookRecordIndex{T0003}$J$-全纯曲线的\BookRecordIndex{T0119}非单射点集至多可数，且只能在 $u$ 的有限个\BookRecordIndex{T0233}临界点处聚集。
发展第 3、6 章的\BookRecordIndex{T0166}横截性理论只需要这些事实。不过，四维 \BookRecordIndex{T0003}$J$-全纯曲线的许多有趣
应用建立在\BookRecordIndex{T0186}交点正性和\BookRecordIndex{T0055}伴随公式之上（见第 9 章）；证明这些结果需要本附录所介绍
的精细局部分析。

设 $J$ 是 $\mathbb C^n$ 上的\BookRecordIndex{T0088}殆复结构，且 $J(0)=J_0$，其中 $J_0:=i$ 为
$\mathbb C^n$ 的标准复结构。记复平面中的闭单位圆盘为
\[
  \mathbb D:=\{z\in\mathbb C\mid |z|\le1\},
\]
并设 $u:(\mathbb D,0)\to(\mathbb C^n,0)$ 是非常值 \BookRecordIndex{T0003}$J$-全纯曲线。考察 $u$
在原点的 Taylor 展开（见引理~\ref{lem:2-4-1}），可知它具有形式
\begin{equation}\label{eq:E-1-1}
  u(z)=az^k+O(|z|^{k+1}),
\end{equation}
其中 $k\ge1$ 为整数，$a\in\mathbb C^n$ 非零。这表明 $u$ 在 $z=0$ 处有
切空间，即直线 $\mathbb Ca$（见注~\ref{rem:E-5-11}）。本附录的结果建立在
Micallef--White~\srcite{287} 给出的远为精细的\BookRecordIndex{T0199}局部标准形之上。该定理断言：
通过源上的 $C^2$-坐标变换和靶上的 $C^1$-坐标变换，$u$ 可共轭到如下形式的
全纯映射：
\[
  v(z)=z^k(a+p(z)),
\]
其中 $a\in\mathbb C^n$ 为非零向量，$p:\mathbb C\to\mathbb C^n$ 为在 $z=0$
处为零、取值于与 $a$ 正交的超平面的多项式。此外，对有限条曲线，可通过靶上的
同一个坐标变换同时得到这种标准形。我们将看到，\BookRecordIndex{T0186}交点正性和\BookRecordIndex{T0055}伴随公式很容易由此\BookRecordIndex{T0326}推出。

\begin{theorem}[Micallef--White]\label{thm:E-1-1}
设 $u:\Sigma\to M$ 为一条 \BookRecordIndex{T0003}$J$-全纯曲线，$\Omega\subset\Sigma$ 为闭包紧的
开集，$x_0\in u(\Omega)\setminus u(\partial\Omega)$，且
\[
  \{z_1,\ldots,z_N\}=u^{-1}(x_0)\cap\Omega.
\]
则存在定义在各 $z_i$ 的两两不交开邻域 $U_i\subset\Omega$ 上的 $C^2$-坐标图
$\phi_i:(U_i,z_i)\to(\mathbb C,0)$、酉矩阵 $L_i\in\mathbb C^{n\times n}$，
以及定义在 $x_0$ 的开邻域 $W\subset M$ 上的 $C^1$-坐标图
$\Psi:(W,x_0)\to(\mathbb C^n,0)$，满足：
\begin{enumerate}
  \item[(i)] $\Psi_*J(x_0)=J_0$，即
    $d\Psi(x_0)J(x_0)=J_0d\Psi(x_0)$；
  \item[(ii)] $u_i:=\Psi\circ u\circ\phi_i^{-1}:\phi_i(U_i)\to\mathbb C^n$
    是变量 $z$ 的多项式；
  \item[(iii)] $L_i u_i(z)=(z^{k_i},z^{k_i}p_i(z))$，其中 $k_i$ 为正整数，
    $p_i:\mathbb C\to\mathbb C^{n-1}$ 为 $z$ 的多项式且 $p_i(0)=0$。
    此外，矩阵 $L_i$ 仅依赖于 $u_i$ 在原点的切空间，即
    \[
      L_i^{-1}(\mathbb C\times\{0\})
      =L_j^{-1}(\mathbb C\times\{0\})\quad\Longrightarrow\quad L_i=L_j.
    \]
\end{enumerate}
\end{theorem}

证明分为三个主要步骤。由于结论是局部的，可将 $(W,x_0)$ 与
$(\mathbb C^n,0)$ 等同，并使 $\mathbb C^n$ 上的\BookRecordIndex{T0088}殆复结构在原点为标准结构。
于是函数 $u\circ\phi_i^{-1}:(\mathbb D,0)\to(\mathbb C^n,0)$ 可按 $z$ 与
$\bar z$ 的幂展开。为控制这些展开，E.5 节考虑一列递减的函数空间
$\mathcal O_\ell\supset\mathcal O_{\ell+1}$，其中 $\mathcal O_\ell$ 由穿孔圆盘
$\mathbb D\setminus\{0\}$ 上的光滑函数 $f$ 组成，它们在原点附近由常数倍
$|z|^\ell$ 控制（精确定义还包括 $f$ 的各阶导数的增长条件）。第一步改进
\eqref{eq:E-1-1}，找出 $C^2$-坐标图 $\phi_i$ 和复矩阵 $L_i$，使复合映射
$u_i:=L_i\circ u\circ\phi_i^{-1}$ 具有 $z\mapsto(z^{k_i},F_i(z))$ 的形式，
其中 $F_i\in\mathcal O_{k_i+1}$。这一步相对直接（见命题~\ref{prop:E-5-10}）。

第二步更为深入：对任意 $i\ne j$，证明差 $F_i-F_j$ 的最低阶项是 $z$ 的某次
幂 $z^\ell$ 的倍数，其中 $\ell>\min(k_i,k_j)$（见命题~\ref{prop:E-6-1}）。
论证在 $F_j=0$ 时也适用，因此还表明 $F_i$ 的最低阶项是全纯的。这一步使用
E.4 节的 \BookRecordIndex{T0024}Hartman--Wintner 定理，它是 Aronszajn~\srcite{18} 的定理在二维
情形下的精化。该定理表明 $F_i-F_j$ 的最低阶项是变量 $s,t$ 的齐次多项式
（其中 $z=s+it$）；由 Cauchy--Riemann 方程的局部分析可知，此多项式是全纯的
（见引理~\ref{lem:E-5-9}）。

最后一步是 Micallef--White 的一个优美论证：利用第二步所得信息，在靶空间 $M$
上构造合适的 $C^1$-坐标图。其思路是“聚焦值域”，使 $F_i-F_j$ 这类原先只近似
等于 $z$ 的幂的函数，精确地等于相应幂。注意，即使在可积情形，前两步是直接的，
最后一步也并非如此。例如，若 $u(z)=(z^k,z^{k+1}F(z))$，其中 $F$ 是任意全纯
函数，如何找到 $C^1$-微分同胚 $\Psi:(\mathbb C^n,0)\to(\mathbb C^n,0)$
使 $\Psi\circ u$ 成为多项式，并不显然。E.7 节实施这一步，且可独立于本附录
其他部分阅读。

定理~\ref{thm:E-1-1} 的第一个应用，是关于单纯 \BookRecordIndex{T0003}$J$-全纯曲线的\BookRecordIndex{T0119}非单射点数目
有限的如下结论。

\begin{theorem}\label{thm:E-1-2}
设 $(M,J)$ 为殆复流形，$(\Sigma,j)$ 为闭黎曼曲面（不要求连通），
$u:\Sigma\to M$ 为单纯 \BookRecordIndex{T0003}$J$-全纯曲线。则 $u$ 的\BookRecordIndex{T0119}非单射点集 $Z\subset\Sigma$ 有限。
\end{theorem}
\begin{proof}
我们证明以下两点。
\begin{enumerate}
  \item[(i)] 对角线在 $\Sigma\times\Sigma$ 中存在开邻域 $U$，使
    \[
      u(z_0)=u(z_1),\quad(z_0,z_1)\in U\quad\Longrightarrow\quad z_0=z_1.
    \]
  \item[(ii)] 自交集
    \[
      S(u):=\{(z_0,z_1)\in\Sigma\times\Sigma\mid
        u(z_0)=u(z_1),\ z_0\ne z_1\}
    \]
    是离散的。
\end{enumerate}
定理立即随之成立：由 (i)，$S(u)$ 是紧空间
$(\Sigma\times\Sigma)\setminus U$ 的闭子集，再由 (ii) 知它有限。
\BookRecordIndex{T0119}非单射点集 $Z$ 是 $S(u)$ 在第一因子上的投影与 $u$ 的\BookRecordIndex{T0233}临界点集之并；
后者由引理~\ref{lem:2-4-1} 知有限，所以 $Z$ 有限。

为证明 (i)，只须证明每个 $z_0\in\Sigma$ 都有邻域 $U_0$，使 $u|_{U_0}$
为单射。一旦如此，用有限多个这样的乘积邻域 $U_j\times U_j$ 覆盖对角线
$\Delta\subset\Sigma\times\Sigma$，并取 $U:=\bigcup_j(U_j\times U_j)$ 即可。
当 $du(z_0)\ne0$ 时，\BookRecordIndex{T0380}隐函数定理立即给出 $u$ 在 $z_0$ 附近的局部单射性。
故设 $du(z_0)=0$。由定理~\ref{thm:E-1-1}，不失一般性可令 $z_0=0$、
$M=\mathbb C^n$，且 $u:\mathbb D\to\mathbb C^n$ 具有形式
\[
  u(z)=(z^k,z^kp(z)),
\]
其中 $k$ 为正整数，$p:\mathbb C\to\mathbb C^{n-1}$ 为在原点为零的多项式。
反设 $u$ 在零附近不局部单射。则存在非零序列 $z_\nu,w_\nu\to0$，使每个 $\nu$
都满足 $u(z_\nu)=u(w_\nu)$ 且 $z_\nu\ne w_\nu$，从而
\[
  z_\nu^k=w_\nu^k,\qquad p(z_\nu)=p(w_\nu).
\]
必要时取子序列，可设存在 $k$ 次单位根 $\zeta\ne1$，使所有 $\nu$ 均有
$w_\nu=\zeta z_\nu$。于是多项式 $z\mapsto p(z)-p(\zeta z)$ 在无穷多个点
$z_\nu$ 处为零，故恒为零。这表明对每个 $z$ 都有 $u(\zeta z)=u(z)$，
与 $u$ 单纯矛盾。这证明了 (i)。

现证 (ii)。取 $(z_0,z_1)\in S(u)$，记 $x:=u(z_0)=u(z_1)$。
由定理~\ref{thm:E-1-1} 可选 $C^1$-坐标图
$\phi_i:(U_i,z_i)\to(\mathbb C,0)$ 与
$\Psi:(W,x)\to(\mathbb C^n,0)$，使映射
$u_i:=\Psi\circ u\circ\phi_i^{-1}$ 具有形式
\[
  u_0(z)=z^k(a+p(z)),\qquad
  u_1(z)=z^\ell(b+q(z)),\qquad p(0)=q(0)=0,
\]
其中 $a,b\in\mathbb C^n$ 非零，$k,\ell$ 为正整数，$p,q:\mathbb C\to
\mathbb C^n$ 为多项式。分两种情形。先设 $a,b$ 线性无关，此时两曲线不相切，
因而交点孤立。更具体地，存在 $\delta>0$，使对所有 $\lambda,\mu\in\mathbb C$，
\[
  |\lambda a+\mu b|\ge\delta(|\lambda|+|\mu|).
\]
又存在常数 $c>0$，使对所有 $z\in\mathbb D$ 有
$|p(z)|\le c|z|$、$|q(z)|\le c|z|$。所以对 $z,w\in\mathbb D$，
\begin{align*}
  |u_0(z)-u_1(w)|
  &\ge |z^ka-w^\ell b|-|z^kp(z)-w^\ell q(w)|\\
  &\ge |z|^k(\delta-c|z|)+|w|^\ell(\delta-c|w|).
\end{align*}
当 $|z|<\delta/c$ 且 $|w|<\delta/c$ 时，右端在 $(z,w)\ne(0,0)$
处为正。这表明，若 $a,b$ 线性无关，则 $(z_0,z_1)$ 是孤立自交点。

现设 $a,b$ 线性相关。由定理~\ref{thm:E-1-1}(iii)，可设
$a=b=(1,0)\in\mathbb C\times\mathbb C^{n-1}$，且 $p(z),q(z)$ 取值于
$\{0\}\times\mathbb C^{n-1}$。反设 $(z_0,z_1)$ 不是 $u$ 的孤立自交点。
则存在非零序列 $z_\nu,w_\nu\to0$，使每个 $\nu$ 都满足
$u_0(z_\nu^\ell)=u_1(w_\nu^k)$，因而
\[
  z_\nu^m=w_\nu^m,\qquad
  p(z_\nu^\ell)=q(w_\nu^k),\qquad m:=k\ell.
\]
必要时取子序列，可设存在 $m$ 次单位根 $\zeta$，使对每个 $\nu$ 有
$w_\nu=\zeta z_\nu$。于是多项式
$z\mapsto p(z^\ell)-q(\zeta^kz^k)$ 在无穷多个点 $z_\nu$ 为零，故恒为零。
因此，对每个 $z\in\mathbb D$ 有
$u_0(z^\ell)=u_1(\zeta^kz^k)$，与 $u$ 单纯矛盾。
这证明了 (ii)，也证明了定理。
\end{proof}

\begin{remark}\label{rem:E-1-3}
定理~\ref{thm:E-1-2} 的证明表明：单纯 \BookRecordIndex{T0003}$J$-全纯曲线 $u:\Sigma\to M$
的每个\BookRecordIndex{T0233}临界点 $z_0\in\Sigma$ 都有一个开邻域 $U$，使 $u|_U$ 单射，
且 $u|_{U\setminus\{z_0\}}$ 为嵌入。这里不要求 $\Sigma$ 为闭曲面。
\end{remark}

\BookRecordIndex{T0036}Micallef--White 定理最重要的推论发生在四维。我们将在 E.2 节证明，在四维情形下
它蕴涵 \BookRecordIndex{T0003}$J$-全纯曲线的\BookRecordIndex{T0186}交点正性。相关结果虽已在 2.6 节陈述，为方便读者，
这里再述一遍。

设 $\Sigma,\Sigma_0,\Sigma_1$ 为闭黎曼曲面，不要求连通；设
$u:\Sigma\to M$、$u_0:\Sigma_0\to M$、$u_1:\Sigma_1\to M$
为 \BookRecordIndex{T0003}$J$-全纯曲线。定义 $u_0$ 与 $u_1$ 的交点数为\sourcefootnote{原书此处定义 $\delta(u_0,u_1)$ 时将条件印作 $u(z_0)=u(z_1)$；由于 $z_i\in\Sigma_i$ 且 $u_i:\Sigma_i\to M$，据映射定义改作 $u_0(z_0)=u_1(z_1)$。}
\[
  \delta(u_0,u_1):=
  \#\{(z_0,z_1)\in\Sigma_0\times\Sigma_1\mid u_0(z_0)=u_1(z_1)\},
\]
并定义 $u$ 的自交点数为
\[
  \delta(u):=\frac12\#\{(z_0,z_1)\in\Sigma\times\Sigma\mid
    u(z_0)=u(z_1),\ z_0\ne z_1\}.
\]
由定理~\ref{thm:E-1-2}，若 $u$ 单纯，则 $\delta(u)$ 有限；若
$u_0,u_1$ 的不交并单纯，则 $\delta(u_0,u_1)$ 有限。

\begin{exercise}\label{ex:E-1-4}
设 $\Sigma_0,\Sigma_1$ 为两个闭连通黎曼曲面，
$u_0:\Sigma_0\to M$、$u_1:\Sigma_1\to M$ 为两条 \BookRecordIndex{T0003}$J$-全纯曲线，
且 $u_0(\Sigma_0)\ne u_1(\Sigma_1)$。证明 $\delta(u_0,u_1)<\infty$。
\end{exercise}

\begin{theorem}[\BookRecordIndex{T0186}交点正性]\label{thm:E-1-5}
设 $(M,J)$ 为四维殆复流形，$A_0,A_1\in H_2(M;\mathbb Z)$ 为分别由单纯
\BookRecordIndex{T0003}$J$-全纯曲线 $u_0:\Sigma_0\to M$ 和 $u_1:\Sigma_1\to M$ 表示的同调类。
假设对任意非空开集 $U_0\subset\Sigma_0$、$U_1\subset\Sigma_1$，
均有 $u_0(U_0)\ne u_1(U_1)$。则
\[
  \delta(u_0,u_1)\le A_0\cdot A_1,
\]
且等号成立当且仅当所有交点都横截，即若
$u_0(z_0)=u_1(z_1)=:x$，则
$\operatorname{im}du_0(z_0)+\operatorname{im}du_1(z_1)=T_xM$。
\end{theorem}

\begin{theorem}[\BookRecordIndex{T0055}伴随公式]\label{thm:E-1-6}
设 $(M,J)$ 为四维殆复流形，$(\Sigma,j)$ 为闭黎曼曲面，不要求连通，
$u:\Sigma\to M$ 为单纯 \BookRecordIndex{T0003}$J$-全纯曲线。记 $A\in H_2(M;\mathbb Z)$
为 $u$ 表示的同调类。则
\begin{equation}\label{eq:E-1-2}
  2\delta(u)\le A\cdot A-c_1(A)+\chi(\Sigma),
\end{equation}
且等号成立当且仅当 $u$ 为浸入且所有自交点横截，即当
$u(z_0)=u(z_1)=:x$ 且 $z_0\ne z_1$ 时，
$\operatorname{im}du(z_0)+\operatorname{im}du(z_1)=T_xM$。
\end{theorem}

\begin{corollary}[McDuff]\label{cor:E-1-7}
设 $M,\Sigma,u,A$ 如定理~\ref{thm:E-1-6}。则
\begin{equation}\label{eq:E-1-3}
  A\cdot A-c_1(A)+\chi(\Sigma)\ge0,
\end{equation}
且等号成立当且仅当 $u$ 为嵌入。
\end{corollary}
\begin{proof}
不等式~\eqref{eq:E-1-3} 由~\eqref{eq:E-1-2} 得出。
\eqref{eq:E-1-3} 中等号成立等价于~\eqref{eq:E-1-2} 中等号成立且
$\delta(u)=0$。
\end{proof}

\begin{remark}\label{rem:E-1-8}
关于~\eqref{eq:E-1-2} 中的亏量
$A\cdot A-c_1(A)+\chi(\Sigma)-2\delta(u)$ 的进一步讨论，见 Milnor~\srcite{290}。
\end{remark}
